\documentclass[10pt,a4paper]{amsart}
\usepackage[margin=19mm]{geometry}
\usepackage{amsmath,amssymb,mathtools}
\usepackage[scr=rsfs,cal=euler]{mathalfa}
\usepackage{xfrac}
\usepackage[unicode,hidelinks,bookmarks=false]{hyperref}
\newcommand{\ie}{i.e.\ }
\newcommand{\cf}{cf.\ }
\newcommand{\resp}{respectively\ }
\newcommand{\Subsection}{\subsection}
\providecommand{\nobreakdash}{\nobreak\hbox{-}}
\setlength{\emergencystretch}{2em}
\multlinegap0pt
\usepackage{bm}
\usepackage{amssymb}
\usepackage{mathtools}
\usepackage[shortlabels]{enumitem}
\usepackage{xcolor}
\newtheorem{lemma}{Lemma}
\newcommand{\D}{\mathcal D}
\newcommand{\E}{\mathcal E}
\newcommand{\Nzero}{\mathbb N_0}
\newcommand{\Vol}{\operatorname{Vol}}
\newcommand{\abs}[1]{\lvert #1\rvert}
\newcommand{\dd}{\,\mathrm d}
\newcommand{\norm}[1]{\lVert #1\rVert}
\newcommand{\wal}{\operatorname{wal}}
\title{Erratum to ``Higher order scrambled digital nets achieve the optimal rate of the root mean square error for smooth integrands''}
\author{Josef Dick}
\date{Source: arXiv:1007.0842v5 (CC BY 4.0)}
\begin{document}
\maketitle

\noindent\emph{Source and licence.} Erratum to ``Higher order scrambled digital nets achieve the optimal rate of the root mean square error for smooth integrands''. Version arXiv:1007.0842v5 (CC BY 4.0), distributed by arXiv under the Creative Commons Attribution 4.0 International licence, http://creativecommons.org/licenses/by/4.0/, as recorded in the arXiv licence metadata for this exact version and shown on the abstract page https://arxiv.org/abs/1007.0842v5. Full licence notice: \texttt{licenses/arxiv-1007.0842-CC-BY-4.0.txt}.\par\smallskip

\noindent\textit{Editorial scope.} Counted unit: the article's lemma lem:block-decay (source0.tex lines 1130--1144). The article's complete original proof of it is delivered with it (source0.tex lines 1146--1836, one \texttt{proof} environment of 691 lines). No mathematics is written, repaired or re-derived: the statement and the proof are the article's own lines, in the article's own notation, and no retained line is substituted or re-flowed.  Retained uncounted as the source-local context the proof needs: lines 1046--1053; lines 1105--1108; lines 1112--1121.  Nothing was omitted from any retained span, so the package carries no omission record.  No retained line contains a citation, so the wrapper carries no bibliography and no bibliography entry.  No retained line contains a figure environment, an image-inclusion command or drawing code, so no image is delivered and the package declares no asset dependency.  Editorial formatting only, in addition: an AMS \texttt{amsart} preamble that restates the article's own declarations and macros used by the retained lines, namely the environments \texttt{lemma} on separate counters, together with the standard packages those lines need.  The retained environments print in this document as Lemma 1 in order of appearance, whereas the article prints the same items with the numbering of its own full document.  The complete native source of the article as distributed in the arXiv e-print for this exact version is bundled unchanged as \texttt{sources/100025/source0.tex}.
\begin{equation}\label{eq:sobolev-norm}
 \norm{f}_{\mathcal H_{s,\alpha}}^2
 :=\abs{I(f)}^2
 +\sum_{\emptyset\ne u\subseteq\{1,\ldots,s\}}
   \ \sum_{\boldsymbol\nu_u\in\{1,\ldots,\alpha\}^{u}}
   \int_{[0,1]^{\abs u}}
   \abs{\partial^{\boldsymbol\nu_u}f_u(x_u)}^2\dd x_u.
\end{equation}

\begin{equation}\label{eq:eta-block}
 \eta(\boldsymbol\ell)
 :=\prod_{i\in v}\prod_{h=1}^{r_i}\eta_{i,h}.
\end{equation}

\begin{lemma}[Order statistics]\label{lem:order-statistics}
Let $0\le a_j\le c_j$ for $1\le j\le n$, and write
$a_{(1)}\le\cdots\le a_{(n)}$ and
$c_{(1)}\le\cdots\le c_{(n)}$ for the corresponding order statistics, with
repetitions.  Then $a_{(h)}\le c_{(h)}$ for every $h$, and therefore
\[
 \prod_{h=1}^r a_{(h)}\le\prod_{h=1}^r c_{(h)}
 \qquad(0\le r\le n).
\]
\end{lemma}

\begin{lemma}[Direct Sobolev replacement for Lemma 9]\label{lem:block-decay}
If $f\in\mathcal H_{s,\alpha}$, then
\begin{equation}\label{eq:sobolev-block-local}
 \sigma_{d,\boldsymbol\ell,s}(f)
 \le 2^{q(\boldsymbol\ell)}b^{\abs K}
      \eta(\boldsymbol\ell)
      \norm{\partial^{\boldsymbol r}f_v}_{L_2([0,1]^{\abs v})},
\end{equation}
where $\boldsymbol r=(r_i)_{i\in v}$.  Consequently,
\begin{equation}\label{eq:sobolev-block}
 \sigma_{d,\boldsymbol\ell,s}(f)
 \le 2^{s\max(d-\alpha,0)}b^{ds}
      \eta(\boldsymbol\ell)\norm f_{\mathcal H_{s,\alpha}}.
\end{equation}
\end{lemma}

\begin{proof}
Write
\[
 L:=|\boldsymbol\ell|_1
   =\sum_{j=1}^{ds}\ell_j,
 \qquad
 K:=\{j\in\{1,\ldots,ds\}:\ell_j>0\}.
\]
For each output coordinate $i\in\{1,\ldots,s\}$, define
\[
 K_i
 :=
 K\cap\{(i-1)d+1,\ldots,id\},
 \qquad
 v:=\{i:K_i\neq\emptyset\},
\]
and put
\[
 r_i:=\min(\alpha,|K_i|),
 \qquad
 q(\boldsymbol\ell)
 :=
 \sum_{i\in v}(|K_i|-r_i).
\]
Thus $r_i$ is the number of increments which will be retained in
output coordinate $i$, while $q(\boldsymbol\ell)$ is the total
number of discarded increments.

We first derive the finite Walsh-block identity which is the starting
point of the argument.

For
$\boldsymbol m=(m_1,\ldots,m_{ds})\in\Nzero^{ds}$, define
\[
 \mathcal A_{\boldsymbol m}
 :=
 \prod_{j=1}^{ds}\{0,\ldots,b^{m_j}-1\}.
\]
In particular,
\[
 \mathcal A_{\boldsymbol\ell}
 =
 \prod_{j=1}^{ds}\{0,\ldots,b^{\ell_j}-1\},
 \qquad
 \mathcal A_1:=\{0,\ldots,b-1\}^{ds}.
\]
For $u\subseteq K$, let $\boldsymbol1_u$ denote its indicator
vector.  Since, for $j\in K$,
\[
 \mathbf 1_{\{b^{\ell_j-1}\le k_j<b^{\ell_j}\}}
 =
 \mathbf 1_{\{0\le k_j<b^{\ell_j}\}}
 -
 \mathbf 1_{\{0\le k_j<b^{\ell_j-1}\}},
\]
multiplication over the active coordinates gives
\begin{equation}\label{eq:block-indicator-detailed}
 \mathbf1_{B_{\boldsymbol\ell}}(\boldsymbol k)
 =
 \sum_{u\subseteq K}
 (-1)^{|u|}
 \mathbf1_{\mathcal A_{\boldsymbol\ell-\boldsymbol1_u}}
 (\boldsymbol k).
\end{equation}
Consequently,
\begin{equation}\label{eq:block-inclusion-exclusion-detailed}
 \sigma_{d,\boldsymbol\ell,s}^2(f)
 =
 \sum_{u\subseteq K}(-1)^{|u|}
 \sum_{\boldsymbol k\in
       \mathcal A_{\boldsymbol\ell-\boldsymbol1_u}}
 \left|
   \widehat f(\E_d(\boldsymbol k))
 \right|^2.
\end{equation}

Equivalently, define the finite Walsh polynomial
\[
 G_{\boldsymbol\ell}(x)
 :=
 \sum_{\boldsymbol k\in B_{\boldsymbol\ell}}
 \widehat f(\E_d(\boldsymbol k))
 \wal_{\E_d(\boldsymbol k)}(x).
\]
By orthogonality of the Walsh system,
\begin{equation}\label{eq:block-parseval-detailed}
 \sigma_{d,\boldsymbol\ell,s}^2(f)
 =
 \int_{[0,1]^s}|G_{\boldsymbol\ell}(x)|^2\,\dd x.
\end{equation}
Using \eqref{eq:block-indicator-detailed}, we may also write
\begin{equation}\label{eq:G-inclusion-exclusion-detailed}
 G_{\boldsymbol\ell}(x)
 =
 \sum_{u\subseteq K}(-1)^{|u|}
 \sum_{\boldsymbol k\in
       \mathcal A_{\boldsymbol\ell-\boldsymbol1_u}}
 \widehat f(\E_d(\boldsymbol k))
 \wal_{\E_d(\boldsymbol k)}(x).
\end{equation}

All frequencies occurring in this expression vanish in the output
coordinates outside $v$.  After integrating out those inactive
coordinates, the corresponding Walsh coefficients may therefore be
computed from $f_v$.  We work from now on on
$[0,1]^{|v|}$.

For $a\in\mathcal A_{\boldsymbol\ell}$, define the fine interlaced
cell
\[
 E_a
 :=
 \operatorname{pr}_v
 \D_d\bigl(
 [ab^{-\boldsymbol\ell},
  (a+1)b^{-\boldsymbol\ell})
 \bigr).
\]
The cells $E_a$, $a\in\mathcal A_{\boldsymbol\ell}$, are pairwise
disjoint up to their boundaries and partition $[0,1]^{|v|}$.  They
fix exactly
\[
 L=\sum_{j\in K}\ell_j
\]
base-$b$ digits, and hence
\begin{equation}\label{eq:Ea-volume-detailed}
 \Vol(E_a)=b^{-L}.
\end{equation}

For $u\subseteq K$, put
\[
 \boldsymbol m_u:=\boldsymbol\ell-\boldsymbol1_u
\]
and define $\pi_u(a)\in\mathcal A_{\boldsymbol m_u}$ by
\[
 [\pi_u(a)]_j
 :=
 \begin{cases}
  \lfloor a_j/b\rfloor, & j\in u,\\
  a_j, & j\notin u.
 \end{cases}
\]
Let
\[
 E_{u,a}
 :=
 \operatorname{pr}_v
 \D_d\bigl(
 [\pi_u(a)b^{-\boldsymbol m_u},
  (\pi_u(a)+1)b^{-\boldsymbol m_u})
 \bigr).
\]
Thus $E_{u,a}$ is obtained from $E_a$ by forgetting the last
prescribed digit in each input coordinate $j\in u$.

We shall use the finite Walsh-kernel identity
\begin{equation}\label{eq:finite-walsh-kernel-detailed}
 \sum_{\boldsymbol k\in\mathcal A_{\boldsymbol m}}
 \wal_{\E_d(\boldsymbol k)}(x)
 \overline{\wal_{\E_d(\boldsymbol k)}(y)}
 =
 b^{|\boldsymbol m|_1}
 \mathbf1_{\{x,y\text{ agree in all digits prescribed by }
                  \boldsymbol m\}}.
\end{equation}
This follows coordinatewise from the one-dimensional identity
\[
 \sum_{k=0}^{b^m-1}
 \wal_k(x)\overline{\wal_k(y)}
 =
 b^m
 \mathbf1_{\{\lfloor b^m x\rfloor=\lfloor b^m y\rfloor\}},
\]
together with the fact that $\E_d$ interlaces the corresponding
frequency digits in the same way that $\D_d$ interlaces the point
digits.

For $x\in E_a$, applying
\eqref{eq:finite-walsh-kernel-detailed} gives
\begin{align}
 &\sum_{\boldsymbol k\in
       \mathcal A_{\boldsymbol\ell-\boldsymbol1_u}}
 \widehat f(\E_d(\boldsymbol k))
 \wal_{\E_d(\boldsymbol k)}(x)
 \notag\\
 &\qquad
 =
 b^{L-|u|}
 \int_{E_{u,a}} f_v(t)\,\dd t.
 \label{eq:lower-projection-cell-detailed}
\end{align}
Hence $G_{\boldsymbol\ell}$ is constant on every $E_a$, and
\begin{equation}\label{eq:G-on-cell-detailed}
 G_{\boldsymbol\ell}(x)
 =
 b^L
 \sum_{u\subseteq K}
 (-1)^{|u|}b^{-|u|}
 \int_{E_{u,a}}f_v(t)\,\dd t,
 \qquad x\in E_a.
\end{equation}
Using \eqref{eq:Ea-volume-detailed} in
\eqref{eq:block-parseval-detailed}, we obtain the exact identity
\begin{equation}\label{eq:block-cell-exact-detailed}
 \sigma_{d,\boldsymbol\ell,s}^2(f)
 =
 b^L
 \sum_{a\in\mathcal A_{\boldsymbol\ell}}
 \left|
   \sum_{u\subseteq K}
   (-1)^{|u|}b^{-|u|}
   \int_{E_{u,a}}f_v(t)\,\dd t
 \right|^2.
\end{equation}

We next resolve each parent cell $E_{u,a}$ into its fine children.
For $a\in\mathcal A_{\boldsymbol\ell}$, write
\[
 e_j(a):=b\lfloor a_j/b\rfloor,
 \qquad
 d_j(a):=a_j-e_j(a)
        =a_j-b\lfloor a_j/b\rfloor.
\]
Thus $d_j(a)\in\{0,\ldots,b-1\}$ is the last base-$b$ digit
of $a_j$.

For $u\subseteq K$ and
$\kappa_u\in\{0,\ldots,b-1\}^{u}$, let
$a^{(u,\kappa)}$ be the fine-child index defined by
\[
 a^{(u,\kappa)}_j
 :=
 \begin{cases}
  e_j(a)+\kappa_j, & j\in u,\\
  a_j, & j\notin u.
 \end{cases}
\]
Then, up to boundaries,
\[
 E_{u,a}
 =
 \bigcup_{\kappa_u\in\{0,\ldots,b-1\}^{u}}
 E_{a^{(u,\kappa)}},
\]
and therefore
\begin{equation}\label{eq:parent-children-detailed}
 \int_{E_{u,a}}f_v
 =
 \sum_{\kappa_u\in\{0,\ldots,b-1\}^{u}}
 \int_{E_{a^{(u,\kappa)}}}f_v.
\end{equation}
Since the summand depends only on the components of the replacement
digit vector in $u$, we may equivalently average over the full set
$\mathcal A_1=\{0,\ldots,b-1\}^{ds}$:
\begin{equation}\label{eq:full-digit-average-detailed}
 b^{-|u|}
 \int_{E_{u,a}}f_v
 =
 b^{-ds}
 \sum_{k\in\mathcal A_1}
 \int_{E_{a^{(u,k)}}}f_v.
\end{equation}
Indeed, each choice of the digits $k_j$, $j\in u$, occurs exactly
$b^{ds-|u|}$ times in the full sum over $k$.

We now express the child cells as translates of $E_a$.
If
\[
 j=(i-1)d+h\in K,
 \qquad 1\le h\le d,
\]
define
\begin{equation}\label{eq:z-aj-detailed}
 z_j(a,k)
 :=
 (k_j-d_j(a))b^{-h-(\ell_j-1)d}.
\end{equation}
Changing the final prescribed digit in input coordinate $j$ from
$d_j(a)$ to $k_j$ changes the corresponding digit of output
coordinate $i$, namely the digit in position
\[
 h+d(\ell_j-1),
\]
and therefore translates that output coordinate by exactly
$z_j(a,k)$.

For $u\subseteq K$, let $Z_u(a,k)\in\mathbb R^{|v|}$ have
$i$th active component
\[
 [Z_u(a,k)]_i
 :=
 \sum_{j\in u\cap K_i}z_j(a,k).
\]
Then, up to boundaries,
\[
 E_{a^{(u,k)}}=E_a+Z_u(a,k),
\]
and translation invariance of Lebesgue measure gives
\begin{equation}\label{eq:child-translate-detailed}
 \int_{E_{a^{(u,k)}}}f_v(x)\,\dd x
 =
 \int_{E_a}f_v(t+Z_u(a,k))\,\dd t.
\end{equation}

Define
\begin{equation}\label{eq:delta-explicit-detailed}
 \delta_{a,k}(t)
 :=
 \sum_{u\subseteq K}(-1)^{|u|}
 f_v\bigl(t+Z_u(a,k)\bigr).
\end{equation}
Combining
\eqref{eq:full-digit-average-detailed} and
\eqref{eq:child-translate-detailed}, we obtain
\begin{align}
 &\sum_{u\subseteq K}
 (-1)^{|u|}b^{-|u|}
 \int_{E_{u,a}}f_v
 \notag\\
 &\qquad
 =
 b^{-ds}
 \sum_{k\in\mathcal A_1}
 \int_{E_a}\delta_{a,k}(t)\,\dd t.
 \label{eq:block-difference-rearrangement-detailed}
\end{align}
Substituting this identity into
\eqref{eq:block-cell-exact-detailed} yields the exact finite
Walsh-block identity
\begin{equation}\label{eq:block-exact-detailed}
 \sigma_{d,\boldsymbol\ell,s}^2(f)
 =
 b^{L-2ds}
 \sum_{a\in\mathcal A_{\boldsymbol\ell}}
 \left|
   \sum_{k\in\mathcal A_1}
   \int_{E_a}\delta_{a,k}(t)\,\dd t
 \right|^2.
\end{equation}
Up to this point all steps are equalities; no Sobolev estimate has
yet been used.

By the triangle inequality,
\[
 \left|
   \sum_{k\in\mathcal A_1}
   \int_{E_a}\delta_{a,k}(t)\,\dd t
 \right|
 \le
 \sum_{k\in\mathcal A_1}
 \int_{E_a}|\delta_{a,k}(t)|\,\dd t.
\]
Set
\[
 A_{a,k}
 :=
 \int_{E_a}|\delta_{a,k}(t)|\,\dd t.
\]
Then
\begin{equation}\label{eq:before-max-detailed}
 \sigma_{d,\boldsymbol\ell,s}^2(f)
 \le
 b^{L-2ds}
 \sum_{a\in\mathcal A_{\boldsymbol\ell}}
 \left(
   \sum_{k\in\mathcal A_1}A_{a,k}
 \right)^2.
\end{equation}

By Cauchy--Schwarz in the finite sum over $k$, using
$|\mathcal A_1|=b^{ds}$,
\[
 \left(
   \sum_{k\in\mathcal A_1}A_{a,k}
 \right)^2
 \le
 b^{ds}\sum_{k\in\mathcal A_1}A_{a,k}^2.
\]
A second application of Cauchy--Schwarz, now on $E_a$, gives
\[
 A_{a,k}^2
 \le
 \Vol(E_a)
 \int_{E_a}|\delta_{a,k}(t)|^2\,\dd t
 =
 b^{-L}
 \int_{E_a}|\delta_{a,k}(t)|^2\,\dd t.
\]
Substituting these estimates into
\eqref{eq:before-max-detailed} gives
\begin{equation}\label{eq:block-L2-reduction-detailed}
 \sigma_{d,\boldsymbol\ell,s}^2(f)
 \le
 b^{-ds}
 \sum_{k\in\mathcal A_1}
 \sum_{a\in\mathcal A_{\boldsymbol\ell}}
 \int_{E_a}|\delta_{a,k}(t)|^2\,\dd t.
\end{equation}

We now estimate the finite differences.  Fix $a$ and $k$.
From \eqref{eq:z-aj-detailed},
\[
 |z_j(a,k)|
 \le
 (b-1)b^{-h-(\ell_j-1)d}
 =:\eta'_j,
 \qquad
 j=(i-1)d+h\in K.
\]
For a fixed output coordinate $i\in v$, arrange the numbers
$|z_j(a,k)|$, $j\in K_i$, in nondecreasing order and retain the
$r_i$ smallest.  Finite-difference operators in the same coordinate
commute, so the order of the increments is irrelevant.

Likewise arrange the numbers $\eta'_j$, $j\in K_i$, in
nondecreasing order.  Since
\[
 |z_j(a,k)|\le\eta'_j,
\]
Lemma~\ref{lem:order-statistics} shows that the $h$th smallest actual
increment is no larger than the $h$th smallest upper bound.  Hence
the product of the $r_i$ retained absolute increments is at most
the product of the $r_i$ smallest $\eta'_j$.  Multiplication over
$i\in v$, together with the definition \eqref{eq:eta-block}, gives
\begin{equation}\label{eq:retained-product-detailed}
 \prod_{i\in v}
 \prod_{\substack{j\in K_i\\ j\ {\rm retained}}}
 |z_j(a,k)|
 \le
 \eta(\boldsymbol\ell).
\end{equation}

We next reduce the order of the finite difference.  Write
\[
 T_yg(t):=g(t+y).
\]
Up to the harmless global sign determined by the convention in
\eqref{eq:delta-explicit-detailed}, the contribution from output
coordinate $i$ is
\[
 \prod_{j\in K_i}(I-T_{z_j e_i}).
\]
Let $R_i\subseteq K_i$ be the indices of the $r_i$ retained
increments and let
\[
 Q_i:=K_i\setminus R_i.
\]
Then
\begin{align}
 \prod_{j\in K_i}(I-T_{z_j e_i})
 &=
 \left(
   \prod_{j\in Q_i}(I-T_{z_j e_i})
 \right)
 \left(
   \prod_{j\in R_i}(I-T_{z_j e_i})
 \right)
 \notag\\
 &=
 \sum_{U_i\subseteq Q_i}
 (-1)^{|U_i|}
 T_{\sum_{j\in U_i}z_j e_i}
 \prod_{j\in R_i}(I-T_{z_j e_i}).
 \label{eq:difference-reduction-detailed}
\end{align}
Thus a difference of order $|K_i|$ is an alternating sum of
$2^{|K_i|-r_i}$ translates of differences of order $r_i$.
Multiplication over all active output coordinates shows that
$\delta_{a,k}$ is an alternating sum of exactly
\[
 \prod_{i\in v}2^{|K_i|-r_i}
 =
 2^{q(\boldsymbol\ell)}
\]
translates of mixed finite differences of order
$\boldsymbol r=(r_i)_{i\in v}$.

For a signed increment $z$, write
\[
 I(z):=[\min(0,z),\max(0,z)].
\]
Repeated use of the fundamental theorem of calculus gives, up to an
irrelevant sign,
\begin{align}
 \Delta_{\boldsymbol r}g(t;\boldsymbol z)
 &=
 \int_{I(z_{1,1})}\!\!\cdots
 \int_{I(z_{i,h})}\!\!\cdots
 \partial^{\boldsymbol r}g
 \left(
   t+\sum_{i,h}\theta_{i,h}e_i
 \right)
 \prod_{i,h}\dd\theta_{i,h}.
 \label{eq:FTC-mixed-detailed}
\end{align}
Consequently, if $S\subseteq[0,1]^{|v|}$ and the translated points
which occur in the integral remain in $[0,1]^{|v|}$, Minkowski's
integral inequality gives
\begin{align}
 \left\|
 \Delta_{\boldsymbol r}g(\,\cdot\,;\boldsymbol z)
 \right\|_{L_2(S)}
 &\le
 \int_{I(z_{1,1})}\!\!\cdots
 \int_{I(z_{i,h})}\!\!\cdots
 \left\|
 \partial^{\boldsymbol r}g
 \left(
   \,\cdot\,
   +\sum_{i,h}\theta_{i,h}e_i
 \right)
 \right\|_{L_2(S)}
 \prod_{i,h}\dd\theta_{i,h}
 \notag\\
 &\le
 \left(\prod_{i,h}|z_{i,h}|\right)
 \left\|
   \partial^{\boldsymbol r}g
 \right\|_{L_2([0,1]^{|v|})}.
 \label{eq:finite-difference-L2-detailed}
\end{align}

It remains to organize the sum over the fine cells $E_a$.  The
increments in \eqref{eq:z-aj-detailed} depend on $a$ only through
the last active digits
\[
 \rho(a):=(d_j(a))_{j\in K}
 \in\{0,\ldots,b-1\}^{K}.
\]
For
\[
 \rho\in\{0,\ldots,b-1\}^{K},
\]
define
\[
 S_\rho
 :=
 \bigcup_{\substack{
 a\in\mathcal A_{\boldsymbol\ell}\\
 d_j(a)=\rho_j\ {\rm for\ all}\ j\in K}}
 E_a.
\]
The union is disjoint up to boundaries.  If $\rho(a)=\rho$, then,
for fixed $k$, all increments $z_j(a,k)$ are the same; denote them
by $z_j(\rho,k)$.  Hence the finite-difference expression
$\delta_{a,k}$ is the same function of its argument for every
$a$ in this group.  Denote it by $\delta_{\rho,k}$.  Then
\begin{equation}\label{eq:group-last-digits-detailed}
 \sum_{\substack{
 a\in\mathcal A_{\boldsymbol\ell}\\
 \rho(a)=\rho}}
 \int_{E_a}|\delta_{a,k}(t)|^2\,\dd t
 =
 \int_{S_\rho}|\delta_{\rho,k}(t)|^2\,\dd t.
\end{equation}

For $t\in S_\rho$, every vertex which occurs in
$\delta_{\rho,k}(t)$ belongs to $[0,1]^{|v|}$.  Indeed, adding
$z_j(\rho,k)$ changes the relevant prescribed interlaced digit
from $\rho_j$ to $k_j$ and leaves all other prescribed digits
unchanged.  The boxes parametrized by the variables
$\theta_{i,h}$ in \eqref{eq:FTC-mixed-detailed} are contained in
the convex hull of the corresponding finite-difference vertices.
Since the cube $[0,1]^{|v|}$ is convex, these boxes remain in the
cube.

We may therefore apply
\eqref{eq:finite-difference-L2-detailed} to each of the
$2^{q(\boldsymbol\ell)}$ reduced differences.  By the triangle
inequality in $L_2$, followed by
\eqref{eq:retained-product-detailed},
\begin{equation}\label{eq:delta-L2-detailed}
 \|\delta_{\rho,k}\|_{L_2(S_\rho)}
 \le
 2^{q(\boldsymbol\ell)}
 \eta(\boldsymbol\ell)
 \left\|
   \partial^{\boldsymbol r}f_v
 \right\|_{L_2([0,1]^{|v|})}.
\end{equation}
Squaring and using
\eqref{eq:group-last-digits-detailed}, we obtain
\[
 \sum_{\substack{
 a\in\mathcal A_{\boldsymbol\ell}\\
 \rho(a)=\rho}}
 \int_{E_a}|\delta_{a,k}(t)|^2\,\dd t
 \le
 4^{q(\boldsymbol\ell)}
 \eta(\boldsymbol\ell)^2
 \left\|
   \partial^{\boldsymbol r}f_v
 \right\|_2^2.
\]

There are exactly $b^{|K|}$ possible vectors
$\rho\in\{0,\ldots,b-1\}^{K}$.  Therefore, for every fixed
$k\in\mathcal A_1$,
\begin{equation}\label{eq:sum-a-bound-detailed}
 \sum_{a\in\mathcal A_{\boldsymbol\ell}}
 \int_{E_a}|\delta_{a,k}(t)|^2\,\dd t
 \le
 b^{|K|}
 4^{q(\boldsymbol\ell)}
 \eta(\boldsymbol\ell)^2
 \left\|
   \partial^{\boldsymbol r}f_v
 \right\|_2^2.
\end{equation}
Substituting \eqref{eq:sum-a-bound-detailed} into
\eqref{eq:block-L2-reduction-detailed} and using
$|\mathcal A_1|=b^{ds}$, we obtain
\begin{align*}
 \sigma_{d,\boldsymbol\ell,s}^2(f)
 &\le
 b^{-ds}b^{ds}b^{|K|}
 4^{q(\boldsymbol\ell)}
 \eta(\boldsymbol\ell)^2
 \left\|
   \partial^{\boldsymbol r}f_v
 \right\|_2^2\\
 &=
 4^{q(\boldsymbol\ell)}b^{|K|}
 \eta(\boldsymbol\ell)^2
 \left\|
   \partial^{\boldsymbol r}f_v
 \right\|_2^2.
\end{align*}
Taking square roots yields
\begin{equation}\label{eq:sobolev-block-local-strong-detailed}
 \sigma_{d,\boldsymbol\ell,s}(f)
 \le
 2^{q(\boldsymbol\ell)}b^{|K|/2}
 \eta(\boldsymbol\ell)
 \left\|
   \partial^{\boldsymbol r}f_v
 \right\|_2.
\end{equation}
Since $b^{|K|/2}\le b^{|K|}$, this implies
\[
 \sigma_{d,\boldsymbol\ell,s}(f)
 \le
 2^{q(\boldsymbol\ell)}b^{|K|}
 \eta(\boldsymbol\ell)
 \left\|
   \partial^{\boldsymbol r}f_v
 \right\|_{L_2([0,1]^{|v|})},
\]
which proves \eqref{eq:sobolev-block-local}.

Finally,
\[
 |K_i|-r_i
 =
 |K_i|-\min(\alpha,|K_i|)
 =
 \max(|K_i|-\alpha,0)
 \le
 \max(d-\alpha,0),
\]
because $|K_i|\le d$.  Summing over at most $s$ output coordinates
gives
\[
 q(\boldsymbol\ell)
 \le
 s\max(d-\alpha,0).
\]
Moreover,
\[
 |K|\le ds.
\]
By the definition of the Sobolev norm in
\eqref{eq:sobolev-norm},
\[
 \left\|
   \partial^{\boldsymbol r}f_v
 \right\|_{L_2([0,1]^{|v|})}
 \le
 \|f\|_{\mathcal H_{s,\alpha}}.
\]
Consequently,
\[
 \sigma_{d,\boldsymbol\ell,s}(f)
 \le
 2^{s\max(d-\alpha,0)}
 b^{ds}
 \eta(\boldsymbol\ell)
 \|f\|_{\mathcal H_{s,\alpha}},
\]
which is \eqref{eq:sobolev-block}.
\end{proof}

\end{document}
